\haddockmoduleheading{Data.Maybe}
\label{module:Data.Maybe}
\haddockbeginheader
{\haddockverb\begin{verbatim}
module Data.Maybe (
    Maybe(Nothing, Just), maybe, isJust, isNothing, fromJust, fromMaybe,
    listToMaybe, maybeToList, catMaybes, mapMaybe
  ) where\end{verbatim}}
\haddockendheader

\section{The Maybe type and operations}
\begin{haddockdesc}
\item[\begin{tabular}{@{}l}
data Maybe a
\end{tabular}]
{\haddockbegindoc
The \haddockid{Maybe} type encapsulates an optional value.  A value of type
 \haddocktt{\haddockid{Maybe} a} either contains a value of type \haddocktt{a} (represented as \haddocktt{\haddockid{Just} a}),
 or it is empty (represented as \haddockid{Nothing}).  Using \haddockid{Maybe} is a good way to
 deal with errors or exceptional cases without resorting to drastic
 measures such as \haddockid{error}.\par
The \haddockid{Maybe} type is also a monad.  It is a simple kind of error
 monad, where all errors are represented by \haddockid{Nothing}.  A richer
 error monad can be built using the \haddockid{Either} type.\par
\enspace \emph{Constructors}\par
\haddockbeginconstrs
\haddockdecltt{=} & \haddockdecltt{Nothing} & \\
\haddockdecltt{|} & \haddockdecltt{Just a} & \\
\end{tabulary}\par}
\end{haddockdesc}
\begin{haddockdesc}
\item[\begin{tabular}{@{}l}
instance MonadZip Maybe\\instance Eq1 Maybe\\instance Ord1 Maybe\\instance Read1 Maybe\\instance Show1 Maybe
\end{tabular}]
\end{haddockdesc}
\begin{haddockdesc}
\item[\begin{tabular}{@{}l}
instance Alternative Maybe
\end{tabular}]
{\haddockbegindoc
Picks the leftmost \haddockid{Just} value, or, alternatively, \haddockid{Nothing}.\par}
\end{haddockdesc}
\begin{haddockdesc}
\item[\begin{tabular}{@{}l}
instance Applicative Maybe\\instance Functor Maybe\\instance Monad Maybe
\end{tabular}]
\end{haddockdesc}
\begin{haddockdesc}
\item[\begin{tabular}{@{}l}
instance MonadPlus Maybe
\end{tabular}]
{\haddockbegindoc
Picks the leftmost \haddockid{Just} value, or, alternatively, \haddockid{Nothing}.\par}
\end{haddockdesc}
\begin{haddockdesc}
\item[\begin{tabular}{@{}l}
instance MonadFail Maybe\\instance Foldable Maybe\\instance Traversable Maybe
\end{tabular}]
\end{haddockdesc}
\begin{haddockdesc}
\item[\begin{tabular}{@{}l}
instance Semigroup a => Monoid Maybe a
\end{tabular}]
{\haddockbegindoc
Lift a semigroup into \haddockid{Maybe} forming a \haddockid{Monoid} according to
 \url{http://en.wikipedia.org/wiki/Monoid}: "Any semigroup \haddocktt{S} may be
 turned into a monoid simply by adjoining an element \haddocktt{e} not in \haddocktt{S}
 and defining \haddocktt{e*e = e} and \haddocktt{e*s = s = s*e} for all \haddocktt{s ∈ S}."\par
\emph{Since 4.11.0}: constraint on inner \haddocktt{a} value generalised from
 \haddockid{Monoid} to \haddockid{Semigroup}.\par}
\end{haddockdesc}
\begin{haddockdesc}
\item[\begin{tabular}{@{}l}
instance Semigroup a => Semigroup Maybe a\\instance Read a => Read Maybe a\\instance Show a => Show Maybe a\\instance Eq a => Eq Maybe a\\instance Ord a => Ord Maybe a
\end{tabular}]
\end{haddockdesc}
\begin{haddockdesc}
\item[\begin{tabular}{@{}l}
maybe :: b -> (a -> b) -> Maybe a -> b
\end{tabular}]
{\haddockbegindoc
The \haddockid{maybe} function takes a default value, a function, and a \haddockid{Maybe}
 value.  If the \haddockid{Maybe} value is \haddockid{Nothing}, the function returns the
 default value.  Otherwise, it applies the function to the value inside
 the \haddockid{Just} and returns the result.\par
\subsubsection*{\textbf{Examples}}
Basic usage:\par
\begin{quote}
{\haddockverb\begin{verbatim}
>>> maybe False odd (Just 3)
True

\end{verbatim}}
\end{quote}
\begin{quote}
{\haddockverb\begin{verbatim}
>>> maybe False odd Nothing
False

\end{verbatim}}
\end{quote}
Read an integer from a string using \haddockid{readMaybe}. If we succeed,
 return twice the integer; that is, apply \haddocktt{(*2)} to it. If instead
 we fail to parse an integer, return \haddocktt{0} by default:\par
\begin{quote}
{\haddockverb\begin{verbatim}
>>> import GHC.Internal.Text.Read ( readMaybe )

>>> maybe 0 (*2) (readMaybe "5")
10

>>> maybe 0 (*2) (readMaybe "")
0

\end{verbatim}}
\end{quote}
Apply \haddockid{show} to a \haddocktt{Maybe Int}. If we have \haddocktt{Just n}, we want to show
 the underlying \haddockid{Int} \haddocktt{n}. But if we have \haddockid{Nothing}, we return the
 empty string instead of (for example) "Nothing":\par
\begin{quote}
{\haddockverb\begin{verbatim}
>>> maybe "" show (Just 5)
"5"

>>> maybe "" show Nothing
""

\end{verbatim}}
\end{quote}}
\end{haddockdesc}
\begin{haddockdesc}
\item[\begin{tabular}{@{}l}
isJust :: Maybe a -> Bool
\end{tabular}]
{\haddockbegindoc
The \haddockid{isJust} function returns \haddockid{True} iff its argument is of the
 form \haddocktt{Just {\char '137}}.\par
\subsubsection*{\textbf{Examples}}
Basic usage:\par
\begin{quote}
{\haddockverb\begin{verbatim}
>>> isJust (Just 3)
True

\end{verbatim}}
\end{quote}
\begin{quote}
{\haddockverb\begin{verbatim}
>>> isJust (Just ())
True

\end{verbatim}}
\end{quote}
\begin{quote}
{\haddockverb\begin{verbatim}
>>> isJust Nothing
False

\end{verbatim}}
\end{quote}
Only the outer constructor is taken into consideration:\par
\begin{quote}
{\haddockverb\begin{verbatim}
>>> isJust (Just Nothing)
True

\end{verbatim}}
\end{quote}}
\end{haddockdesc}
\begin{haddockdesc}
\item[\begin{tabular}{@{}l}
isNothing :: Maybe a -> Bool
\end{tabular}]
{\haddockbegindoc
The \haddockid{isNothing} function returns \haddockid{True} iff its argument is \haddockid{Nothing}.\par
\subsubsection*{\textbf{Examples}}
Basic usage:\par
\begin{quote}
{\haddockverb\begin{verbatim}
>>> isNothing (Just 3)
False

\end{verbatim}}
\end{quote}
\begin{quote}
{\haddockverb\begin{verbatim}
>>> isNothing (Just ())
False

\end{verbatim}}
\end{quote}
\begin{quote}
{\haddockverb\begin{verbatim}
>>> isNothing Nothing
True

\end{verbatim}}
\end{quote}
Only the outer constructor is taken into consideration:\par
\begin{quote}
{\haddockverb\begin{verbatim}
>>> isNothing (Just Nothing)
False

\end{verbatim}}
\end{quote}}
\end{haddockdesc}
\begin{haddockdesc}
\item[\begin{tabular}{@{}l}
fromJust :: HasCallStack => Maybe a -> a
\end{tabular}]
{\haddockbegindoc
The \haddockid{fromJust} function extracts the element out of a \haddockid{Just} and
 throws an error if its argument is \haddockid{Nothing}.\par
\subsubsection*{\textbf{Examples}}
Basic usage:\par
\begin{quote}
{\haddockverb\begin{verbatim}
>>> fromJust (Just 1)
1

\end{verbatim}}
\end{quote}
\begin{quote}
{\haddockverb\begin{verbatim}
>>> 2 * (fromJust (Just 10))
20

\end{verbatim}}
\end{quote}
\begin{quote}
{\haddockverb\begin{verbatim}
>>> 2 * (fromJust Nothing)
*** Exception: Maybe.fromJust: Nothing
...

\end{verbatim}}
\end{quote}
WARNING: This function is partial. You can use case-matching instead.\par}
\end{haddockdesc}
\begin{haddockdesc}
\item[\begin{tabular}{@{}l}
fromMaybe :: a -> Maybe a -> a
\end{tabular}]
{\haddockbegindoc
The \haddockid{fromMaybe} function takes a default value and a \haddockid{Maybe}
 value.  If the \haddockid{Maybe} is \haddockid{Nothing}, it returns the default value;
 otherwise, it returns the value contained in the \haddockid{Maybe}.\par
\subsubsection*{\textbf{Examples}}
Basic usage:\par
\begin{quote}
{\haddockverb\begin{verbatim}
>>> fromMaybe "" (Just "Hello, World!")
"Hello, World!"

\end{verbatim}}
\end{quote}
\begin{quote}
{\haddockverb\begin{verbatim}
>>> fromMaybe "" Nothing
""

\end{verbatim}}
\end{quote}
Read an integer from a string using \haddockid{readMaybe}. If we fail to
 parse an integer, we want to return \haddocktt{0} by default:\par
\begin{quote}
{\haddockverb\begin{verbatim}
>>> import GHC.Internal.Text.Read ( readMaybe )

>>> fromMaybe 0 (readMaybe "5")
5

>>> fromMaybe 0 (readMaybe "")
0

\end{verbatim}}
\end{quote}}
\end{haddockdesc}
\begin{haddockdesc}
\item[\begin{tabular}{@{}l}
listToMaybe :: {\char 91}a{\char 93} -> Maybe a
\end{tabular}]
{\haddockbegindoc
The \haddockid{listToMaybe} function returns \haddockid{Nothing} on an empty list
 or \haddocktt{\haddockid{Just} a} where \haddocktt{a} is the first element of the list.\par
\subsubsection*{\textbf{Examples}}
Basic usage:\par
\begin{quote}
{\haddockverb\begin{verbatim}
>>> listToMaybe []
Nothing

\end{verbatim}}
\end{quote}
\begin{quote}
{\haddockverb\begin{verbatim}
>>> listToMaybe [9]
Just 9

\end{verbatim}}
\end{quote}
\begin{quote}
{\haddockverb\begin{verbatim}
>>> listToMaybe [1,2,3]
Just 1

\end{verbatim}}
\end{quote}
Composing \haddockid{maybeToList} with \haddockid{listToMaybe} should be the identity
 on singleton/empty lists:\par
\begin{quote}
{\haddockverb\begin{verbatim}
>>> maybeToList $ listToMaybe [5]
[5]

>>> maybeToList $ listToMaybe []
[]

\end{verbatim}}
\end{quote}
But not on lists with more than one element:\par
\begin{quote}
{\haddockverb\begin{verbatim}
>>> maybeToList $ listToMaybe [1,2,3]
[1]

\end{verbatim}}
\end{quote}}
\end{haddockdesc}
\begin{haddockdesc}
\item[\begin{tabular}{@{}l}
maybeToList :: Maybe a -> {\char 91}a{\char 93}
\end{tabular}]
{\haddockbegindoc
The \haddockid{maybeToList} function returns an empty list when given
 \haddockid{Nothing} or a singleton list when given \haddockid{Just}.\par
\subsubsection*{\textbf{Examples}}
Basic usage:\par
\begin{quote}
{\haddockverb\begin{verbatim}
>>> maybeToList (Just 7)
[7]

\end{verbatim}}
\end{quote}
\begin{quote}
{\haddockverb\begin{verbatim}
>>> maybeToList Nothing
[]

\end{verbatim}}
\end{quote}
One can use \haddockid{maybeToList} to avoid pattern matching when combined
 with a function that (safely) works on lists:\par
\begin{quote}
{\haddockverb\begin{verbatim}
>>> import GHC.Internal.Text.Read ( readMaybe )

>>> sum $ maybeToList (readMaybe "3")
3

>>> sum $ maybeToList (readMaybe "")
0

\end{verbatim}}
\end{quote}}
\end{haddockdesc}
\begin{haddockdesc}
\item[\begin{tabular}{@{}l}
catMaybes :: {\char 91}Maybe a{\char 93} -> {\char 91}a{\char 93}
\end{tabular}]
{\haddockbegindoc
The \haddockid{catMaybes} function takes a list of \haddockid{Maybe}s and returns
 a list of all the \haddockid{Just} values.\par
\subsubsection*{\textbf{Examples}}
Basic usage:\par
\begin{quote}
{\haddockverb\begin{verbatim}
>>> catMaybes [Just 1, Nothing, Just 3]
[1,3]

\end{verbatim}}
\end{quote}
When constructing a list of \haddockid{Maybe} values, \haddockid{catMaybes} can be used
 to return all of the "success" results (if the list is the result
 of a \haddockid{map}, then \haddockid{mapMaybe} would be more appropriate):\par
\begin{quote}
{\haddockverb\begin{verbatim}
>>> import GHC.Internal.Text.Read ( readMaybe )

>>> [readMaybe x :: Maybe Int | x <- ["1", "Foo", "3"] ]
[Just 1,Nothing,Just 3]

>>> catMaybes $ [readMaybe x :: Maybe Int | x <- ["1", "Foo", "3"] ]
[1,3]

\end{verbatim}}
\end{quote}}
\end{haddockdesc}
\begin{haddockdesc}
\item[\begin{tabular}{@{}l}
mapMaybe :: (a -> Maybe b) -> {\char 91}a{\char 93} -> {\char 91}b{\char 93}
\end{tabular}]
{\haddockbegindoc
The \haddockid{mapMaybe} function is a version of \haddockid{map} which can throw
 out elements.  In particular, the functional argument returns
 something of type \haddocktt{\haddockid{Maybe} b}.  If this is \haddockid{Nothing}, no element
 is added on to the result list.  If it is \haddocktt{\haddockid{Just} b}, then \haddocktt{b} is
 included in the result list.\par
\subsubsection*{\textbf{Examples}}
Using \haddocktt{\haddockid{mapMaybe} f x} is a shortcut for \haddocktt{\haddockid{catMaybes} {\char '44} \haddockid{map} f x}
 in most cases:\par
\begin{quote}
{\haddockverb\begin{verbatim}
>>> import GHC.Internal.Text.Read ( readMaybe )

>>> let readMaybeInt = readMaybe :: String -> Maybe Int

>>> mapMaybe readMaybeInt ["1", "Foo", "3"]
[1,3]

>>> catMaybes $ map readMaybeInt ["1", "Foo", "3"]
[1,3]

\end{verbatim}}
\end{quote}
If we map the \haddockid{Just} constructor, the entire list should be returned:\par
\begin{quote}
{\haddockverb\begin{verbatim}
>>> mapMaybe Just [1,2,3]
[1,2,3]

\end{verbatim}}
\end{quote}}
\end{haddockdesc}